\section{Scanning collars and improving the entropy bound}
\label{sec:scanner}

The preceding estimates compare the entropy gained by moving a small set
across a partition with the energy cost of changing its metric.  We now
construct a finite schedule of such moves.  Random initial offsets make
any fixed interaction unlikely to cross a moving boundary; random choices
of local moves give each crossing interaction a chance to be moved in full.
Together these two features produce a small-energy interpolation.
The rough norm comparison then improves the safe-box exponent, and the
sharp comparison gives a sublinear mutual-information bound across a collar.

Fix the cut \(A\), and write \(Z\) for its crossing-edge endpoints.
Throughout this section, depth means ambient integer sup-norm distance.
Interaction supports and their truncations still use the graph metric of
\(\Lambda\).  In particular, a depth row is a set of sites at a fixed
distance from the target, and need not be a horizontal lattice row.

\subsection{The moving partitions}

Let \(T\subset\Z^2\) be a finite ambient target.  Write
\(T_j=T+[-j,j]^2_{\Z}\), \(Q_j=A\cap(T_j\setminus T)\),
and \(X=A\cap T\), and suppose \(X\ne\varnothing\).
A useful example is a safe integer rectangle of maximal side length
\(s_0\): depth rows in a collar of width at most \(s_0\) contain
\(O(s_0)\) sites, so the row scale can be a fixed multiple of \(s_0\).
The safe-box estimate bounds its core entropy; covering a sufficiently
narrow shell by safe squares gives the shell entropy input.  We give
this verification, and its extension to the polygonal unions needed by
the final tiling, after deriving the general scanner estimate.

At an integer scale \(n\ge2\), choose
\begin{equation}
 L=\lfloor n^{1-\ell}\rfloor,
 \quad m=\lfloor n^\mu\rfloor,
 \quad K=\lfloor L/(8m)\rfloor,
 \quad D=\lceil n^\kappa\rceil,
 \quad a=W/K,
 \label{scanner:scales}
\end{equation}
where \(0<\ell<1\), \(0<\kappa<\mu<1-\ell\), and \(W\ge1\).  We work at large
enough \(n\) that \(K\ge1\).  Assume
\begin{equation}
 |T|\le Cn^2,\qquad
 |T_j\setminus T_{j-1}|\le n\quad(1\le j\le L),\qquad
 \dist_\infty(T,Z)>2L+10r_0,
 \label{scanner:test-geometry}
\end{equation}
where \(r_0=\lceil C\log^2 n\rceil\) is chosen as in
Proposition~\ref{prop:truncation}.  Increasing its constant fixes the
truncation error at \(n^{-1000}\).  We require \(n\) large enough that
\(r_0\ll D\ll m\ll L\).  The compact set used for that proposition
is \(S_0=A\cap T_L\); it has \(O(n^2)\) sites.  Denote the resulting
Hamiltonian, exact ground vector, ground energy, and spectral gap by
\(\widetilde H,\widetilde\Omega,\widetilde E_0\), and
\(\widetilde g\), respectively.  Write \(\widetilde h_i\) for its
terms and \(\widetilde X_i\) for their designated supports, as in
that proposition.  In particular
\(\widetilde E_0\le n^{-1000}\) and \(\widetilde g\ge g/2\).
The target cut has \(O(n)\) edges: an ambient edge leaving \(T\)
has its outside endpoint in the first positive-depth row, and the
clearance excludes edges from \(X\) to \(\Lambda\setminus A\).
This observation will control the target's marginal spectrum even for
these general targets.

\begin{center}
\begin{tabular}{@{}ll@{}}
\toprule
Parameter & Role in the scan\\
\midrule
$n$ & Number of padded slots in a depth row\\
$L$ & Total collar depth\\
$m$ & Window scale and number of possible offsets\\
$K$ & Number of separated bands\\
$D$ & Charge lookahead in depth\\
$r_0$ & Graph radius of every splitting support\\
$a,\ W=aK$ & Per-band and total scalar metric weights\\
\bottomrule
\end{tabular}
\end{center}
The scalar $W$ is distinct from the subsystem operators $W_{Q,k}$.

Order the sites of each positive-depth row, and pad that order to exactly
\(n\) slots.  An absent site or a site outside \(A\) is a blank for
the scan.  The order of these slots will remain fixed.  For band
\(g'=0,\ldots,K-1\), put \(b_{g'}=8g'm\), and choose an independent
uniform offset \(r_{g'}\in\{0,\ldots,m-1\}\).  Its initial physical
partition is
\begin{align*}
 P_0&=\{x\in A:\dist_\infty(x,T)\le b_{g'}+m+r_{g'}\}=XU,\\
 V&=(\Lambda\setminus A)\cup
       \{x\in A:\dist_\infty(x,T)>b_{g'}+5m+r_{g'}\},\\
 Y&=\Lambda\setminus(P_0\cup V).
\end{align*}
The two auxiliary systems of Proposition~\ref{prop:comparators}, denoted
here by \(C\) and \(R\), give the full partition
\(P=CP_0\), \(Y\), \(F=RV\).  Every band uses the same auxiliary
systems, but has its own evolving physical partition.

There are two kinds of rounds.  A \emph{fill} advances a nominal front
according to a fixed schedule.  A \emph{charge} chooses an interaction
and attempts to move the entire unassigned part of its designated support.
The random choices specify the weights of a finite history tree; the
metric construction below uses the whole tree.

There are \(nm\) fill rounds and \(nm\) charge rounds, interleaved,
with a fill first.  A round performs one transition in every band.
At fill rounds the receiving side alternates between the near and far
sides.  Its schedule consumes one slot of the original middle window,
in increasing depth for the near side and decreasing depth for the far
side.  Within a row use its fixed order, reversed on the far side if
necessary.  If the slot is a site still in \(Y\), move it to the
receiving side.  Consume the slot also when it is blank or was already
assigned.  In particular, a charge never changes the fill schedule.

For the near side use oriented depth \(\lambda=\dist_\infty(\cdot,T)\),
and for the far side use \(\lambda=-\dist_\infty(\cdot,T)\).  If
\(t_P,t_F\) are the numbers of fills already assigned to the respective
sides, their next incomplete row depths are
\begin{equation}
 j_P=b_{g'}+m+r_{g'}+1+\lfloor t_P/n\rfloor,
 \qquad
 j_F=-(b_{g'}+5m+r_{g'})+\lfloor t_F/n\rfloor.
 \label{scanner:fronts}
\end{equation}
At a charge round choose a uniform side.  On that side form the list of
labelled term anchors with oriented depth in
\([j-r_0,j+D]\), counting anchor multiplicity, and pad it to a fixed
integer \(M_n=\lceil C_1nD\rceil\).  Here \(C_1\) is a uniform
constant large enough for every such list.  Choose a uniform slot of
this list.  If its graph \(r_0\)-ball meets both the receiving physical
side and \(Y\), move all sites of that ball which remain in \(Y\)
to the side.  Otherwise make no move.  The choices are conditionally
independent across bands.  The band margins ensure that every unsigned
depth appearing in these lists is between 1 and \(L\).

The two probabilities needed later have different roles.  Initial offsets
will bound the probability that a fixed term splits by \(O(D/m)\).
Except on rare histories, a charge will select each term currently split
with probability at least \(c/(nD)\).  The deterministic fill schedule
is essential to both estimates: a charge never changes a nominal front.
Figure~\ref{scanner:figure} shows the nominal partitions and the possible
lead of the charged sites.

A history records the offsets and all choices already made.  Its
probability is the classical path probability, before transporting any
replica state.  At each round form the conditional binary mean trees
of the new choices and
interpolate every old leaf by the same parameter \(p\in[0,1]\), as in
Proposition~\ref{prop:transport}.  At \(p=1\) attach the choice trees
to the history tree and proceed to the next round.  The parameter \(p\)
is common to all old histories and all bands in that round.  This gives
the root metric \(\mathcal M\), the norm \(N\), and the normalized
replica vector \(v\) used in the preceding two sections.

\begin{figure}[H]
\centering
\begin{tikzpicture}[x=.83cm,y=1cm,font=\small,>=Stealth]
 \foreach \y/\leftfront/\rightfront/\band in
   {1.25/3.2/6.2/{g},0/8.1/11.1/{g+1}}{
  \fill[blue!15] (0,\y) rectangle (.9,{\y+.65});
  \fill[green!12] (.9,\y) rectangle (\leftfront,{\y+.65});
  \fill[yellow!22] (\leftfront,\y) rectangle (\rightfront,{\y+.65});
  \fill[gray!15] (\rightfront,\y) rectangle (13,{\y+.65});
  \draw (0,\y) rectangle (13,{\y+.65});
  \draw (.9,\y)--(.9,{\y+.65});
  \draw[densely dashed] (\leftfront,{\y-.08})--(\leftfront,{\y+.83});
  \draw[densely dashed] (\rightfront,{\y-.08})--(\rightfront,{\y+.83});
  \fill[pattern=north east lines,pattern color=gray!65]
       (\leftfront,\y) rectangle ({\leftfront+.55},{\y+.65});
  \fill[pattern=north east lines,pattern color=gray!65]
       ({\rightfront-.55},\y) rectangle (\rightfront,{\y+.65});
  \node at (.45,{\y+.325}) {$X$};
  \node at ({(.9+\leftfront)/2},{\y+.325}) {$U$};
  \node at ({(\leftfront+\rightfront)/2},{\y+.325}) {$Y$};
  \node at ({(\rightfront+13)/2},{\y+.325}) {$V$};
  \node[anchor=east] at (-.15,{\y+.325}) {band $\band$};
  \draw[->] (\leftfront,{\y+.8})--({\leftfront+.9},{\y+.8});
  \draw[->] (\rightfront,{\y+.8})--({\rightfront-.9},{\y+.8});
 }
 \draw[<->] (3.2,2.35)--node[above] {$D+O(r_0)$}(3.75,2.35);
 \draw[->] (0,-.5)--(13,-.5) node[right] {depth};
\end{tikzpicture}
\caption{Schematic nominal depth partitions in two bands, not to scale.
The horizontal coordinate records distance from the target, not a physical
lattice direction.  Dashed lines are the nominal fronts, and arrows show
the fill directions.  Charges may have assigned individual sites in the
hatched intervals, so the actual sets \(U,Y,V\) need not be depth
intervals.  The separated windows ensure that one support cannot split
two bands.  Every site of \(\Lambda\setminus A\) also belongs to \(V\)
in each band.  Auxiliary systems are not shown.}
\label{scanner:figure}
\end{figure}

\begin{lemma}[Geometry and probabilities of the scan]
\label{scanner:histories}
Under \eqref{scanner:scales}--\eqref{scanner:test-geometry}, the preceding
construction has the following properties for all sufficiently large \(n\).
\begin{enumerate}
\item Every splitting designated support is a graph \(r_0\)-ball wholly
inside \(A\).  It splits at most one band, and there it meets \(Y\)
and exactly one of the two sides.  The metric algebras generated by all
statuses of distinct bands commute on the symmetric replica space.
\item Every assigned site of \(A\) has oriented depth at most \(j+D+r_0\) on
its side.  Each actual set \(U\) or \(UY\) is comparable by inclusion
with a deterministic positive-depth prefix, and differs from that prefix
by at most \(CnD\) sites.  The cuts of \(U,UY,Y,P_0,V\) have at most
\(CnD\) edges.  At most \(CKnD\) terms split in any joint status.
\item For every labelled term \(i\), every round, and every common
\(p\), the sum of the full terminal-leaf probabilities on which it
splits is at most \(CD/m\).
\item An old history is called good if no side has an assigned site of \(A\) of
oriented depth greater than \(j+D/2\).  For every fixed \(M>0\), the
probability of a history failing this condition at any time is at most
\(n^{-M}\), for sufficiently large \(n\).  At a charge round on a good
old history, each split term is sampled on its splitting side with
probability at least \(c/(nD)\).
\end{enumerate}
All probabilities here are the classical path probabilities of the
history tree, before the transport of replica states.
\end{lemma}

\begin{proof}
\textbf{Deterministic geometry and nesting.}
A truncated support meeting \(S_0\) has radius \(r_0\): when its
anchor is at distance \(d>2r_0\), the alternative radius
\(\lfloor d/2\rfloor\) cannot reach \(S_0\).  Every splitting
support meets a compact part contained in \(S_0\), so this observation
applies.  If an \(r_0\)-ball meeting \(S_0\) also contained a site of
\(\Lambda\setminus A\), a graph path inside the ball would contain
a cut edge with an endpoint at ambient distance at most \(L+2r_0\)
from \(T\).  This contradicts \eqref{scanner:test-geometry}.

At a charge the largest newly assigned oriented depth is at most
\(j+D+r_0\); subsequent movement of the nominal front only improves
this bound.  Each side receives at most \(\lceil nm/2\rceil\) fills,
so advances through at most \(m/2+O(1)\) deterministic rows.  Its
charged lead is \(D+O(r_0)=o(m)\).  The two fronts in a window remain
separated by order \(m\), as do the active regions of successive
windows.  This proves both that an opposite side cannot have consumed
a scheduled slot and that a support cannot split two fronts.  All nonblank slots
strictly behind a nominal front have been assigned to that side.

A near set contains its completed deterministic prefix and is contained
in the prefix ending at its maximal allowed lead.  The difference has
at most \(CnD\) sites.  For \(UY\), take the complement of the
assigned far suffix within the original shell prefix: it is contained
in the deterministic prefix ending at the nominal far front and differs
from it only in \(O(D)\) rows.  These deterministic comparison prefixes
are all prefixes of one fixed ordered list of positive-depth slots and
are consequently nested, even across histories and bands.

A full depth prefix has \(O(n)\) boundary edges.  For its outer
boundary at positive depth, use the last included row, since depth
changes by at most one along an edge.  For the boundary at depth zero,
use the first positive row.  Partial rows cost at most \(4n\) more
edges.  There are no compact-to-opposite-color edges, by the clearance.
Changing at most \(CnD\) sites then gives the claimed \(O(nD)\)
bound for all status cuts and their differences.  Split anchors themselves
lie in \(O(D)\) rows around the nominal fronts, containing \(O(nD)\)
anchors per band.  Summing gives \(CKnD\).

For two bands \(g'<g''\), the entire near side together with the middle
of the first band is contained in the near side of the second band,
for every pair of histories.  After adjoining \(C\), this says
\(P_{g'}Y_{g'}\subset P_{g''}\).  Mirror the central function on
\(F_{g'}\) to its complement \(P_{g'}Y_{g'}\), using
Lemma~\ref{lem:schur}.  The resulting functions from the first band
commute with those of the second: they are disjoint from its \(Y,F\)
parts and invariant under the joint permutations whose center acts on
its \(P\) part.  This proves the asserted cross-band commutation.

\smallskip
\noindent\textbf{Dilution by the initial offsets.}
For a fixed anchor to split, its oriented depth must lie between
\(j-O(r_0)\) and \(j+D+O(r_0)\).  Both formulas in
\eqref{scanner:fronts} are affine, with coefficient \(1\) or \(-1\),
in the initial uniform offset.  Thus only \(O(D)\) of the \(m\)
offsets can permit this event, and only \(O(1)\) bands can reach the
fixed anchor.  This proves \(CD/m\) separately for the old and new
leaf distributions, and hence also for their common-\(p\) mixture.
It does not require conditional independence of the split event and
the later charge choices.

\smallskip
\noindent\textbf{Charge ancestry and availability.}
We finally bound the probability that charged sites get too far ahead of
a nominal front.  This will ensure that every old split lies in the
charge list.  Suppose a side has an assigned site of \(A\) leading
its nominal front by more than
\(D/2\).  Trace the site's first assignment backwards through the
charge ball that assigned it, choosing a previously assigned same-side
site in that ball, until an initial or deterministic assignment is
reached.  An ancestry with \(j'\) charges changes oriented depth by
at most \(2r_0j'\).  If its first charge is at time \(t_1\) and the
history is tested at time \(t\), then
\[
 2r_0j'>D/2+j(t)-j(t_1)-O(1).
\]
The fixed fill schedule gives
\(j(t)-j(t_1)\ge (t-t_1)/(4n)-O(1)\).  Substitution in the
preceding inequality yields both
\[
 j'\ge cD/r_0,\qquad t-t_1\le Cnr_0j'.
\]
Thus all ancestry times lie in this last interval of rounds.
All sites in the ancestry are in \(A\), including for the far side,
because the charging balls are wholly in \(A\).

For a fixed endpoint, test time, and length \(j'\), the time choices
and spatial choices have at most
\[
 \binom{\lceil Cnr_0j'\rceil}{j'}(Cr_0^4)^{j'}
\]
possibilities.  The spatial bound permits both an anchor and a
predecessor within \(O(r_0)\) graph distance at each step, and includes
the bounded anchor multiplicity.  Conditional on earlier choices, a
prescribed hit has probability at most \((cnD)^{-1}\).  The bound
\(\binom{v}{j'}\le(ev/j')^{j'}\) therefore makes the total probability
at most
\[
 \poly(n)\sum_{j'\ge cD/r_0}(Cr_0^5/D)^{j'}.
\]
The endpoints being counted are charged sites in \(S_0\), of which
there are \(O(n^2)\); the numbers of test times and bands are also
polynomial in \(n\).  Thus no far-domain volume enters this factor.
Since \(r_0\) is polylogarithmic and \(D=n^{\kappa+o(1)}\), the bound
is smaller than
every fixed inverse power of \(n\).

On a good old history, a split anchor has oriented depth at least
\(j-r_0\), because its ball meets an unassigned site, and at most
\(j+D/2+r_0\), because it meets the side.  For large \(n\) this
interval lies in the sampling interval \([j-r_0,j+D]\).  Choosing
the relevant side and its labelled slot has probability
\(1/(2M_n)\ge c/(nD)\), and the resulting move is exactly the
\(Y\)-part of that designated support.
\end{proof}

\subsection{Selecting a low-energy interpolation}

The scan provides the nested prefixes required by the norm comparisons
and the two probability estimates required by entropy and energy transport.
We now combine them to select one low-energy interpolation.  The entropy
input below concerns the original ground state; the polynomially accurate
truncation transfers it to \(\widetilde\Omega\).  No gap for a Hamiltonian
restricted to the target or its collar is used.

\begin{proposition}[The scanner estimate]
\label{prop:scanner}
Assume \eqref{scanner:scales} and \eqref{scanner:test-geometry}, and let
\(0<e<1\) satisfy \(\kappa\le e\).  Suppose
\begin{equation}
 S_\Omega(X)\le C_e n^{1+e},\qquad
 S_\Omega(Q_j)\le C_e(nL^e+n)\quad(0\le j\le L).
 \label{scanner:entropy-input}
\end{equation}
Let \(\epsilon=n^{-\nu}\), where \(\nu>0\) is fixed, and suppose
\(a r_0^2\) is sufficiently small.  Construct the auxiliary systems,
the fixed label sector, and the unnormalized vector \(\mathrm{pre}\)
as in Proposition~\ref{prop:comparators}, using typical width
\(w=n^{3/5}\).  For each sufficiently large replica count \(k\),
there are a charge round and a common parameter
\(p_k\in[\epsilon/2,\epsilon]\) at which the normalized vector
\(v_k=\mathcal M^{-s}\mathrm{pre}/N\), \(s=1/4\), has defect energy
\[
 E_{{\rm def},k}
 =\frac{\langle v_k,\overline{\widetilde H}v_k\rangle
                         -\widetilde E_0}{\widetilde g}
\]
bounded by
\begin{equation}
 E_{{\rm def},k}
 \le Cn^\ell D^2W^2(\log n)^C
        \bigl(\delta_n^{1/8}+\epsilon+n^{-100}\bigr)
        +Cn^{-1000}+o_k(1),
 \label{scanner:energy-output}
\end{equation}
where
\begin{equation}
 \delta_n=C_e\epsilon^{-1}
        \bigl(n^{e-\mu}+a^{1/4}(\log n)^C\bigr).
 \label{scanner:density-output}
\end{equation}
Here \(\overline{\widetilde H}\) is the mean of the physical
Hamiltonian over the copies.  All norm comparisons of
Proposition~\ref{prop:comparators} apply at the selected points, with
shell entropy bound \(B_{\rm sh}\le C_e(nL^e+n)\), mismatch
\(B_{\rm exc}\le CnD\), and marginal parameter
\(\mathcal B\le CnD(\log n)^C\).  The sharp comparisons additionally
require their stated condition on \(a\sqrt{\mathcal B}\).

The entire finite physical system, \(n\), the scan, and all auxiliary
dimensions are fixed before \(k\to\infty\).  The remainder is uniform
over the selected round and parameter, although these may depend on \(k\).
Constants may depend on the fixed exponents and the input constant \(C_e\).
\end{proposition}

\begin{proof}
The transfer from \(\Omega\) to \(\widetilde\Omega\) costs
\(O(n^{-248})\) in any entropy in \eqref{scanner:entropy-input}:
the trace distance is \(O(n^{-500})\), and the compact systems have
\(O(n^2)\) sites, so Lemma~\ref{lem:continuity} applies.  A partial
last depth row adds at most \(n\log q\).  Lemma~\ref{scanner:histories}
then gives the matched-prefix and mismatch bounds required by the
comparators.  It also gives the \(O(nD)\) status boundaries.  A
crossing support has Hilbert dimension at most \(q^{C(1+r_0)^2}\),
and its anchor lies within distance \(r_0\) of a cut edge.
Consequently the marginal-tail parameter of Lemma~\ref{lem:tail} is
at most \(CnD(\log n)^C\).  At the target itself it is at most
\(Cn(\log n)^C\).  The latter bound shows that the width \(w=n^{3/5}\)
has failure mass at most \(n^{-100}\) for all sufficiently large
\(n\).  In fact its tail is bounded by
\(C\exp[-c n^{1/10}/(\log n)^C]\).

At a charge round write \(h\) for an old joint history, of probability
\(w_h\), and \(\sigma_{h,u}\) for its transported state in
Proposition~\ref{prop:transport}.  Let \(\mu_\sigma\) denote the
coherent probability measure of Lemma~\ref{lem:symbol}.  The entropy
derivative uses the old-history weights \(w_h\); the same transported
states occur in the energy estimate with terminal weights \((1-p)w_h\).
We therefore measure the entropy cost of the old splits by
\(\mathcal Q(p)\), defined as
\begin{equation}
 \sum_{h\ {\mathrm{good}}}w_h
  \int_{\R}m_s(u)
   \int\sum_{i\ {\mathrm{split\ in}}\ h}
       \eta_{i,h}(\theta)\,d\mu_{\sigma_{h,u}}(\theta)\,du.
 \label{scanner:Q-definition}
\end{equation}
Each split term in this sum has a unique splitting band, and its
partition \(P,Y,F\) is understood to be that band's partition.
If term \(i\) splits on the near side, its moved set is
\(x=Y\cap\widetilde X_i\), the unassigned part of its designated
\(r_0\)-ball, and
\(\eta_{i,h}=S_\theta(x\mid P)+S_\theta(x\mid Y\setminus x)\);
on the far side exchange \(P,F\).  These quantities are nonnegative.
They are also bounded by \(C(\log n)^C\), because the moved sites lie
in an \(r_0\)-ball.  The measure \(m_s(u)\,du\) has fixed finite
mass \(2s\).

On every good old history each split term is sampled with probability
at least \(c/(nD)\), and its sampled move has exactly the corresponding
\(\eta_{i,h}\).  Equation~\eqref{transport:entropy-gain} therefore
gives, at a charge round,
\begin{equation}
 -\partial_p\log N^2
 \ge \frac{cka}{nD}\mathcal Q(p)
       -CkaK a^{1/4}(\log n)^C-O_{\rm fixed}(\log(k+2)).
 \label{scanner:charge-gain}
\end{equation}
At fill rounds retain just the same negative error; the entropy gains
are nonnegative.  Here and below ``fixed'' means that the complete
single-copy instance and the scanner have been fixed.  No such
constant survives the replica limit.

Integrate over every full round, including fills.  Adjacent endpoint
metrics agree by construction, so the left sides telescope.  Initially
\(-\log N^2\ge-O_{\rm fixed}(\log(k+2))\): the floor
\eqref{replicas:leaf-floor} passes to the initial mean tree by
Lemma~\ref{transport:means}, and \(\|\mathrm{pre}\|\le1\).
At the last endpoint the rough
upper comparison \eqref{comparator:rough-upper} gives
\[
 -k^{-1}\log N^2
 \le 2\log d_*-\log z+C_eW(nL^e+nD+n)+o_k(1)
 \le C_eW n^{1+e}+o_k(1).
\]
For the last inequality use \(\kappa\le e\), \(L\le n\), and
\(\log d_*=S_{\widetilde\Omega}(X)+O(w+|\log z|)\).
Restrict the positive charge integrals in
\eqref{scanner:charge-gain} to \([\epsilon/2,\epsilon]\), but pay
the negative errors on all \(2nm\) full rounds.  Since \(aK=W\),
the integrated inequality, normalized by the number of charge rounds, is
\[
 \frac1{nm}\sum_{\text{charge rounds}}
  \int_{\epsilon/2}^{\epsilon}\frac{\mathcal Q(p)}{KnD}\,dp
 \le C_e\left(\frac{n^e}{m}+a^{1/4}(\log n)^C\right)+o_k(1).
\]
The interval has length \(\epsilon/2\), and \(m\asymp n^\mu\).
Averaging therefore gives a charge round and a \(p_k\) in this interval
such that
\begin{equation}
 \frac{\mathcal Q(p_k)}{KnD}\le\delta_n+o_k(1).
 \label{scanner:selected-density}
\end{equation}
The fixed factor two in the interval length is absorbed into constants.
All integrands are bounded at fixed \(k\), and continuous in the
interior interpolation parameter, so ordinary averaging suffices for
this selection.

Apply the energy estimate \eqref{transport:energy} at this same point.
For each term its split weight is at most \(CD/m\).  There are at
most \(CKnD\) split term occurrences in each joint leaf.  On old
good leaves the transported states are exactly those in
\eqref{scanner:Q-definition}; their full terminal weight differs only
by the factor \(1-p_k\le1\).  H\"older's inequality, applied to the
positive measure over histories, terms, \(u\), and \(\theta\), bounds
their sum of \(\eta^{1/8}\) by
\[
 CKnD\bigl(\delta_n+o_k(1)\bigr)^{1/8}.
\]
On bad old histories use the uniform local dimension bound and
Lemma~\ref{scanner:histories} with an arbitrarily large fixed tail
power.  On all new leaves use the same dimension bound and their
total terminal weight \(p_k\le\epsilon\).  After increasing the fixed
logarithmic power, the energy estimate is therefore
\[
 \langle v_k,\overline{\widetilde H}v_k\rangle
 \le 2\widetilde E_0+
 C a^2\frac Dm KnD(\log n)^C
       \bigl(\delta_n^{1/8}+\epsilon+n^{-100}\bigr)+o_k(1).
\]
The factors in this expression have distinct origins: \(D/m\) is the
split probability of one fixed term, \(KnD\) counts split occurrences,
and \(a^2\) is the local skew cost.  Their product simplifies to
\[
 a^2\frac Dm KnD
   =\frac{W^2nD^2}{mK}
   \le CW^2n^\ell D^2.
\]
Subtracting \(\widetilde E_0\) and dividing by
\(\widetilde g\ge g/2\) proves
\eqref{scanner:energy-output}.

Only splitting terms contribute to this energy error; all others have
zero skew exactly.  Thus its coefficient contains no total-volume
factor.  The symbol remainders can be summed over all terms before
taking \(k\to\infty\), because that complete term set is finite and
fixed.  Their uniformity in replica states and in the interpolation
parameter also justifies the possibly \(k\)-dependent selection in
\eqref{scanner:selected-density}.
\end{proof}

\subsection{Geometric targets and entropy inputs}

We now supply the scanner's geometric and entropy inputs for unions of
rectangles and triangles.  These include the geometric targets of the final
tiling, whose number of pieces may grow with the scale.  The constant
\(D_0\) and the function \(F_{\rm box}\) are those of
Proposition~\ref{prop:initial-box}.  Fix the numerical constant
\(C_{\rm tpl}\) large enough to absorb the layer-count constant in the
proof below.  The entropy bounds will be reused each time the safe-box
exponent improves.

\begin{definition}[A template]
\label{scanner:template}
Let \(s_0\ge1\) and \(n\) be positive integers.  A template with parameters
\((n,s_0)\) is a nonempty finite set \(T\subset\Z^2\) of the form
\[
 T=\bigcup_{a=1}^{N_0}(P_a\cap\Z^2),
 \qquad n\ge C_{\rm tpl}N_0(s_0+1),
 \qquad \diam_\infty T\le n,
\]
where \(N_0\ge1\) is an integer and each \(P_a\) is a closed convex rectangle or triangle of sup-norm
diameter at most \(s_0\), all of whose sides have slope
\(0,\infty,1\), or \(-1\).  Empty sampled pieces can be omitted.
For the fixed cut \(A\), the template is separated from the cut if
\[
 \dist_\infty(T,Z)>4D_0s_0.
\]
As above, \(T_j\) is the ambient integer dilation,
\(X=A\cap T\), and \(Q_j=A\cap(T_j\setminus T)\), for \(j\ge0\).
\end{definition}

The diameter condition merely fixes a convenient common size parameter.
The estimates below use the number and sizes of the pieces, rather than
convexity of their union.  They therefore allow the disconnected fragments
which occur in the later tiling.

\begin{lemma}[Rows, coverings, and shell entropy]
\label{scanner:templates}
For \(C_{\rm tpl}\) sufficiently large, a template satisfies
\begin{equation}
 |T|\le Cns_0,\qquad
 |T_j\setminus T_{j-1}|\le n\quad(1\le j\le s_0),
 \qquad |\partial_{\Z^2}T_j|\le Cn\quad(0\le j\le s_0).
 \label{scanner:template-geometry}
\end{equation}
Suppose it is separated from the cut and, for some \(0<e<1\),
\(F_{\rm box}(r)\le C_e r^{1+e}\) for every integer \(r\ge1\).
Then, for every integer \(1\le L\le s_0\),
\begin{equation}
 S_\Omega(X)\le C'_e ns_0^e,\qquad
 S_\Omega(Q_j)\le C'_e nL^e\quad(0\le j\le L).
 \label{scanner:template-entropy}
\end{equation}
The same shell bound, with an additional \(Cn\), holds for any prefix
obtained by including all depth rows before one row and an arbitrary
subset of that last row.  All physical cuts of \(X,Q_j\), and
\(A\cap T_j\) have \(O(n)\) edges.
\end{lemma}

\begin{proof}
Consider first one sampled piece.  On an integer horizontal row its sites
form an integer interval.  Its left endpoint is a maximum of affine
integer-valued functions with slopes in \(\{-1,0,1\}\), and its right
endpoint is a minimum of such functions; horizontal sides impose an
interval constraint on the row index.  To obtain these formulas, round
the intercepts of the defining line inequalities.  Because the slopes
are integers, this rounding is independent of the integer row index.
The difference of the lower and upper bounds is convex, so the nonempty
rows are consecutive.  On consecutive nonempty rows each endpoint moves
by at most one.  Consecutive row intervals overlap or are adjacent as
integer intervals.

After integer sup-norm dilation by \(j\), the interval on a row is the
union of the original intervals in the window of rows at distance at
most \(j\), each extended by \(j\) on both ends.  This union is still
an integer interval.  Its endpoints are the sliding minimum and maximum
of the old endpoints, minus or plus \(j\).  They again move by at most
one between consecutive nonempty rows: the nonempty index windows at
two consecutive rows have Hausdorff distance at most one, and the old
endpoint functions are 1-Lipschitz.  One more dilation changes each
endpoint on an old row by at most two and adds at most two rows, each
of length \(O(s_0+j+1)\).  This proves a bound
\(C(s_0+j+1)\) for both the new layer and the edge boundary of one
piece.  Summing over the pieces proves
\eqref{scanner:template-geometry}, with the layer constant absorbed in
\(C_{\rm tpl}\).  The area bound follows from
\(|P_a\cap\Z^2|\le C(s_0+1)^2\).

We also need a count at intermediate dyadic scales.  In the usual dyadic
partition of the lattice into integer squares, the number of side-\(u\)
squares containing both a site of a dilated piece and a site outside it
is
\begin{equation}
 O\bigl((s_0+j+1)/u+1\bigr),\qquad u\ge1.
 \label{scanner:mixed-piece}
\end{equation}
Indeed, in a block of \(u\) consecutive nonempty rows the two endpoints
vary by at most \(u\).  Only constantly many square columns near each
endpoint can be mixed.  At the first and last blocks count instead the
whole width, which costs \(O((s_0+j+1)/u+1)\) in total.
A square mixed for a union, or for a difference of two unions, must be
mixed for at least one of the constituent sets.  Thus the corresponding
count for \(T\), \(T_j\), or \(T_j\setminus T\), with \(j\le s_0\),
is at most
\(C(N_0+n/u)\).

Cover \(T_j\setminus T\) by maximal contained dyadic squares, with side
capped at the largest power of two not exceeding \(L\).  These squares
partition the set: a single site is an admissible square, and the dyadic
nesting makes distinct maximal squares disjoint.  At the cap there are
at most \(Cn/L\) squares because
\(|T_j\setminus T|\le nj\le nL\).  Below the cap, the parent of each
maximal square is mixed.  Equation~\eqref{scanner:mixed-piece} therefore
bounds the number of squares of side \(u\) by
\(C(N_0+n/u)\le Cn/u\), since \(u\le L\le s_0\) and
\(N_0s_0\le n\).  All these squares are safe: their sites lie in
\(T_j\), their sides are at most \(s_0\), and
\[
 \dist_\infty(T_j,Z)\ge\dist_\infty(T,Z)-j
   >4D_0s_0-s_0>D_0s_0.
\]
Subadditivity and the assumed box bound now give
\[
 S_\Omega(Q_j)
 \le C_e n\sum_{\substack{u\le L\\u\ \mathrm{dyadic}}}u^e
 \le C'_e nL^e.
\]
Covering \(T\) in the same way, with cap comparable to \(s_0\),
uses \(O(n/s_0)\) squares at the cap and \(O(n/u)\) below it, and
gives the first bound in \eqref{scanner:template-entropy}.

Adding an arbitrary part of one depth row changes entropy by at most
\(n\log q\).  Finally, near these sets there is no edge joining \(A\)
to its complement: such an edge would have an endpoint in \(Z\) within
distance \(s_0+1\) of \(T\).  Thus their physical boundaries are bounded
by the ambient template and shell boundaries just counted.  This argument
does not count missing lattice edges as physical boundary edges.
\end{proof}

\subsection{A small safe-box exponent}

We use the scanner first with a bounded total metric weight.  The rough
comparison is sufficient here, even when \(a\sqrt{\mathcal B}\) is
large.  This separates the bootstrap from the sharper concentration
condition needed only in the final application.

\begin{proposition}[Improved safe-box estimate]
\label{prop:small-box}
There is a constant \(C\), depending only on the Hamiltonian parameters,
such that
\[
 F_{\rm box}(r)\le Cr^{1+e_*}\quad(r\ge1),
 \qquad e_*=2\cdot10^{-6}.
\]
\end{proposition}

\begin{proof}
Let \(e_0\in(0,1)\) be supplied by
Proposition~\ref{prop:initial-box}.  If \(e_0\le e_*\), the assertion
already follows.  Otherwise fix, once for the entire bootstrap,
\begin{equation}
 g_0=(1-e_0)/2,
 \quad \mu=1-g_0,
 \quad \nu=g_0/1000,
 \quad \ell=\nu/100,
 \quad \kappa=\min(\nu/100,e_*/4),
 \quad W=4.
 \label{scanner:bootstrap-parameters}
\end{equation}
These parameters will not change when the current entropy exponent
improves.

Suppose \(F_{\rm box}(r)\le C_e r^{1+e}\), where
\(e_*<e\le e_0\).  Take any safe integer rectangle \(T\) of actual
maximal side length \(s_0\), and set \(n=\lceil C_2s_0\rceil\)
for a sufficiently large fixed \(C_2\).  If \(A\cap T\) is empty,
its entropy is zero, so assume otherwise.  For large \(s_0\),
\(L=o(s_0)\), so \eqref{scanner:test-geometry} holds.  In particular
the weaker safe-box clearance \(D_0s_0\), rather than the stronger
template clearance of Definition~\ref{scanner:template}, is enough
for this application.

Here is the shell entropy input explicitly.  The shell
\(T_j\setminus T\), \(j\le L\), has area \(O(nL)\) and consists
of a bounded number of rectangular strips.  Its mixed dyadic squares
of side \(u\) number \(O(n/u+1)\), by the boundary-line count.
The maximal-contained-square argument of
Lemma~\ref{scanner:templates}, capped at side comparable to \(L\),
therefore gives \(S_\Omega(Q_j)\le C_e nL^e\).  These covering
squares are safe: their sides are at most \(L\), and their distance
from \(Z\) is at least \(D_0s_0-L>D_0L\) once \(s_0\) is large.
This also covers thin rectangles.  The target entropy is at most
\(C_e s_0^{1+e}\).  Adding a partial row costs \(O(n)\), so all
the hypotheses of Proposition~\ref{prop:scanner} hold.

We check its quantitative output.  Since
\(K\asymp n^{g_0-\ell}\),
\(a\asymp n^{-g_0+\ell}\).  In
\eqref{scanner:density-output} the two powers before logarithms are
at most
\[
 e_0-\mu+\nu=-0.999g_0,
 \qquad -(g_0-\ell)/4+\nu=-0.2489975g_0.
\]
Thus \(\delta_n\le n^{-0.24g_0}\) for sufficiently large \(n\),
after constants and fixed logarithmic powers have been absorbed.
The energy prefactor has power
\(\ell+2\kappa\le0.00003g_0\).  Its slower decaying input is
\(\epsilon=n^{-0.001g_0}\), rather than
\(\delta_n^{1/8}\le n^{-0.03g_0}\).  Hence
\begin{equation}
 E_{{\rm def},k}\le Cn^{-\nu/2}+o_k(1).
 \label{scanner:bootstrap-energy}
\end{equation}
All scale conditions and the smallness of \(ar_0^2\) hold at large
\(n\).  No sharp marginal-moment condition has been used.

Put \(\tau=n^{-\nu/4}\).  Apply the rough upper and lower
comparisons \eqref{comparator:rough-upper} and
\eqref{comparator:rough-lower} at the selected interpolation.  Write
\(S_X=S_{\widetilde\Omega}(X)\).  The upper pin cost has leading
coefficient \(2S_X\); the lower coefficient is
\(4sWS_X=4S_X\).  Replacing \(\log d_*\) and \(S'\) by \(S_X\)
costs \(O(w+|\log z|)\).  The remaining relative loss in this
coefficient is at most
\(C(\tau+E_{{\rm def},k}/\tau)\), which is small after taking
\(k\to\infty\) and then \(n\) sufficiently large.  The two
comparisons consequently imply
\begin{equation}
 S_X\le C_e\bigl(nL^e+nD+w+h(\tau)/a+1\bigr).
 \label{scanner:bootstrap-comparison}
\end{equation}
To justify the order of limits in this absorption, keep \(n\) fixed
in both norm comparisons and use
\eqref{scanner:bootstrap-energy}.  The resulting scalar inequality
has the same \(S_X\) for every \(k\); all other remainders are
uniform at the selected points and vanish.  Only then choose the
large-scale threshold making the coefficient loss smaller than the
fixed margin between four and two.

For \(0<\tau<1/2\), \(h(\tau)\le C\tau\log(e/\tau)\).  Therefore
\[
 h(\tau)/a
 \le C n^{g_0-\ell-\nu/4}\log n=O(n),
\]
because \(g_0<1/2\).  Also \(w=n^{3/5}\le n\), and
\(\kappa\le e_*/4<(1-\ell)e\) while \(e>e_*\).
Equation~\eqref{scanner:bootstrap-comparison}, followed by the
negligible transfer back to \(\Omega\), gives
\[
 S_\Omega(A\cap T)\le C'_e n^{1+(1-\ell)e}
                       \le C''_e s_0^{1+(1-\ell)e}.
\]
Bounded \(s_0\) are included by increasing the constant, using the
volume bound.  Taking the supremum over safe boxes thus replaces
the current exponent \(e\) by \((1-\ell)e\).

Starting from \(e_0\), finitely many repetitions reach an exponent
at most \(e_*\), which can be loosened to \(e_*\).  Every iteration
uses only the current safe-box estimate on the shell covering squares.
The parameters in \eqref{scanner:bootstrap-parameters}, and the finite
number of iterations, depend only on the fixed Hamiltonian parameters.
Thus the final constant is uniform over the domain and cut.
\end{proof}

\subsection{The collar mutual information}

After the bootstrap we allow the total metric weight to grow slowly.
The exact norm comparison then charges the remaining target pin cost by
the reciprocal of that weight.  This is the point at which the entropy
identity in Proposition~\ref{prop:comparators} gives a bound independent
of the selected history.

\begin{proposition}[A sublinear collar bound]
\label{prop:collar}
Let \(T\) be a template separated from the cut, with parameters
\((n,s_0)\) as in Definition~\ref{scanner:template}.  Set
\[
 L=\lfloor n^{1-10^{-5}}\rfloor,
 \qquad X=A\cap T,
 \qquad S_0=A\cap T_L,
\]
and suppose \(L\le s_0\).  For all sufficiently large \(n\),
\begin{equation}
 I_\Omega(X:\Lambda\setminus S_0)
       \le C n^{1-\varepsilon},\qquad \varepsilon=10^{-5}.
 \label{scanner:collar-output}
\end{equation}
The constant and threshold are independent of \(T,s_0,n,\Lambda\),
and \(A\).  The cuts of \(X,S_0\), and \(S_0\setminus X\) have
\(O(n)\) edges, and \(|S_0|=O(n^2)\).
\end{proposition}

\begin{proof}
If \(X\) is empty there is nothing to prove.  Otherwise
Proposition~\ref{prop:small-box} and
Lemma~\ref{scanner:templates} supply
\eqref{scanner:entropy-input} with \(e=e_*=2\cdot10^{-6}\).
Since \(s_0\le n\), the core entropy is at most \(Cn^{1+e}\).
The template clearance and \(L\le s_0\) imply
\eqref{scanner:test-geometry} at large \(n\); the logarithmic
truncation radius is negligible compared with \(L\).

Use the following fixed parameter choices:
\begin{equation}
 \mu=\frac14,\quad \nu=\frac1{1000},\quad
 \ell=\frac1{100000},\quad
 \omega=\frac1{20000},\quad
 \kappa=\frac1{1000000},\quad W=n^\omega.
 \label{scanner:final-parameters}
\end{equation}
Then \(a\asymp n^{-0.74994}\).  The two terms in
\eqref{scanner:density-output} have powers
\(-0.248998\) and \(-0.186485\), before the fixed logarithmic
factors.  Hence \(\delta_n\le n^{-0.18}\) at sufficiently large
scales.  The power of the energy prefactor is
\[
 \ell+2\kappa+2\omega=0.000112.
\]
The \(\epsilon=n^{-0.001}\) term is again the limiting one, and
Proposition~\ref{prop:scanner} gives
\begin{equation}
 E_{{\rm def},k}\le Cn^{-0.0005}+o_k(1).
 \label{scanner:final-energy}
\end{equation}
Set \(\tau=n^{-0.00025}\).  All fixed multiples of \(a\) lie in
the marginal-moment interval required for the sharp comparisons,
since
\[
 a\sqrt{\mathcal B}
 \le C n^{-0.2499395}(\log n)^C\longrightarrow0.
\]

At each leaf write, as before, the physical partition as \(X,U,Y,V\).
All entropies in the next calculation refer to \(\widetilde\Omega\).
Let
\[
 g_{\rm pre}=S(XU)+S(V)-S(Y),\qquad
 g_{\rm post}=S(U)+S(UY)-S(Y).
\]
These are the quantities in the sharp norm comparisons.  Their upper
bounds are uniform in the status.  In particular
\(S(XU)\le S(X)+S(U)\le Cn^{1+e}\), and by purity
\(S(V)=S(XUY)\le S(X)+S(UY)\le Cn^{1+e}\).
Here the matched-prefix error \(CnD\) is absorbed because
\(\kappa<e\).  Thus \(0\le g_{\rm pre}\le Cn^{1+e}\).

Denote the full terminal path and uniform band average by
\(\langle\cdot\rangle_{\rm leaf}\).  Subtracting
\eqref{comparator:sharp-upper} from
\eqref{comparator:sharp-lower}, and then dividing by \(2sW\), bounds
\(2S(X)+\langle g_{\rm pre}-g_{\rm post}\rangle_{\rm leaf}\)
by a constant times
\begin{equation}
 \frac{S(X)+1}{W}+w
 +(\tau+E_{{\rm def},k}/\tau)n^{1+e}
 +\frac{h(\tau)}a+a\mathcal B
 +\frac{|\log z|}a+o_k(1).
 \label{scanner:final-errors}
\end{equation}
This calculation uses \(S'=S(X)+O(w+|\log z|)\) and the analogous
estimate for \(\log d_*\).  It also restores the lost
\(\tau\langle g_{\rm pre}\rangle_{\rm leaf}\) from the lower
comparison using the preceding uniform bound.

For each leaf separately, purity gives the exact identity
\begin{equation}
 2S(X)+g_{\rm pre}-g_{\rm post}
       =I(X:V)+I(X:YV).
 \label{scanner:entropy-identity}
\end{equation}
Both \(V\) and \(YV\) contain the fixed subsystem
\(\Lambda\setminus S_0\).  Monotonicity of mutual information
therefore makes the left side of \eqref{scanner:entropy-identity}
at least \(I_{\widetilde\Omega}(X:\Lambda\setminus S_0)\),
independently of the history.  This remains true after the leaf average,
even though the selected round and its weights can vary with \(k\).

For completeness, after \(k\to\infty\) the terms in
\eqref{scanner:final-errors} have the following upper bounds on their
powers of \(n\), apart from fixed logarithms.
\begin{center}
\begin{tabular}{@{}lc@{}}
\toprule
Error term & Power of $n$\\
\midrule
$(S(X)+1)/W$ & $0.999952$\\
$w$ & $0.6$\\
$(\tau+E_{{\rm def},k}/\tau)n^{1+e}$, after $k\to\infty$ & $0.999752$\\
$h(\tau)/a$ & $0.74969$\\
$a\mathcal B$ & $0.250061$\\
$|\log z|/a$ & $-99.25006$\\
\bottomrule
\end{tabular}
\end{center}
All are strictly below \(1-\varepsilon=0.99999\).
Equation~\eqref{scanner:final-energy} controls the mass-loss term
uniformly before the limit.  This proves
\(I_{\widetilde\Omega}(X:\Lambda\setminus S_0)
\le Cn^{1-\varepsilon}\).

Finally write mutual information as
\(S(X)-S(X\mid\Lambda\setminus S_0)\) and use
Lemma~\ref{lem:continuity}.  Only \(\log\dim\mathcal H_X=O(n^2)\)
enters, even if the far subsystem is very large.  The trace distance
\(O(n^{-500})\) between the original and truncated ground states
therefore changes this mutual information by \(O(n^{-248})\).
This proves \eqref{scanner:collar-output}.  The final size and boundary
assertions were established in Lemma~\ref{scanner:templates}.
\end{proof}

The replica construction has now been eliminated: the conclusion of
Proposition~\ref{prop:collar} is an inequality for the original finite
ground state.  The next section applies the exact positive quasi-local
system to that state.  In particular, the exponentially large auxiliary
operator used there never multiplies a truncation or a replica error
from the present argument.
